% Unique CQ stimuli relocated from theory topics; no answers or answer-writing guides.
\subsection{বিষয়ভিত্তিক CQ অনুশীলন}
\par\Needspace{7\baselineskip}
\textbf{প্রয়োগ 1: বিষয় ২.3}\par
\par
কলেজে তিন ধরনের যোগাযোগ প্রয়োজন: ল্যাবের কম্পিউটারগুলো নেটওয়ার্কে যুক্ত করা, পুরো ক্যাম্পাসে তারবিহীন ইন্টারনেট সেবা দেওয়া এবং অনলাইন ক্লাসের ভিডিও সার্ভারকে উচ্চগতির ব্যাকবোনে যুক্ত করা।

\begin{enumerate}
\item ক. ট্রান্সমিশন মাধ্যম কী?
\item খ. নির্দেশিত ও অনির্দেশিত মাধ্যমের মূল পার্থক্য কী?
\item গ. ল্যাবের কম্পিউটারগুলো যুক্ত করতে কোন মাধ্যম উপযোগী? ব্যাখ্যা করো।
\item ঘ. ভিডিও সার্ভারের উচ্চগতির ব্যাকবোন সংযোগে অপটিক্যাল ফাইবার অধিক উপযোগী হতে পারে কেন? বিশ্লেষণ করো।
\end{enumerate}
\par

\par\Needspace{7\baselineskip}
\textbf{প্রয়োগ 2: বিষয় ২.3}\par
একটি হাসপাতাল রোগী ব্যবস্থাপনা ডিজিটাল পদ্ধতিতে পরিচালনা করতে চায়। চিকিৎসকদের কক্ষে ডেস্কটপ কম্পিউটার ও রোগীদের অপেক্ষার স্থানে Wi-Fi থাকবে। সার্ভার কক্ষ থেকে ইন্টারনেট ব্যাকবোনে উচ্চগতির সংযোগ প্রয়োজন। দূরবর্তী এলাকার জরুরি সেবায় স্যাটেলাইট সংযোগ বিবেচনা করা হচ্ছে।

\begin{enumerate}
\item ক. নির্দেশিত মাধ্যম কী?
\item খ. বেতার মাধ্যম চলাচলের সুবিধা দেয় কেন?
\item গ. চিকিৎসকদের কক্ষের ডেস্কটপ কম্পিউটারের নেটওয়ার্কে কোন মাধ্যম উপযোগী? ব্যাখ্যা করো।
\item ঘ. সার্ভার কক্ষের উচ্চগতির সংযোগ ও দূরাঞ্চলের সেবায় কোন মাধ্যমগুলো উপযোগী? বিশ্লেষণ করো।
\end{enumerate}
\par

\par\Needspace{7\baselineskip}
\textbf{প্রয়োগ 3: বিষয় ২.4}\par
\par
কলেজে নতুন ICT ল্যাব তৈরি হচ্ছে। ল্যাবের ৫০টি কম্পিউটার সুইচে যুক্ত হবে। অধ্যক্ষের কার্যালয় থেকে সার্ভার কক্ষে উচ্চগতির সংযোগ প্রয়োজন। পাশে উচ্চক্ষমতার বৈদ্যুতিক যন্ত্র থাকায় তড়িৎচৌম্বক বিঘ্নের আশঙ্কা আছে।

\begin{enumerate}
\item ক. নির্দেশিত মাধ্যম কী?
\item খ. টুইস্টেড পেয়ার কেবলের তার পেঁচানো থাকে কেন?
\item গ. ল্যাবের কম্পিউটার যুক্ত করতে কোন কেবল উপযোগী? ব্যাখ্যা করো।
\item ঘ. অধ্যক্ষের কার্যালয় থেকে সার্ভার কক্ষে উচ্চগতির ও বিঘ্নমুক্ত সংযোগে কোন কেবল অধিক উপযোগী? যুক্তিসহ বিশ্লেষণ করো।
\end{enumerate}
\par

\par\Needspace{7\baselineskip}
\textbf{প্রয়োগ 4: বিষয় ২.4}\par
হাসপাতালের নতুন নেটওয়ার্কে চিকিৎসকদের কক্ষের কম্পিউটারগুলো সুইচে যুক্ত হবে। সার্ভার কক্ষ থেকে প্রধান ইন্টারনেট ব্যাকবোনে উচ্চগতির সংযোগ প্রয়োজন। কয়েকটি স্থানে বৈদ্যুতিক যন্ত্র থাকায় তড়িৎচৌম্বক বিঘ্ন হতে পারে।

\begin{enumerate}
\item ক. STP কী?
\item খ. অপটিক্যাল ফাইবারের আলোকপথ তড়িৎচৌম্বক বিঘ্নে প্রভাবিত হয় না কেন?
\item গ. চিকিৎসকদের কক্ষের কম্পিউটার সংযোগে কোন কেবল উপযোগী? ব্যাখ্যা করো।
\item ঘ. সার্ভার কক্ষের উচ্চগতির ব্যাকবোন সংযোগে অপটিক্যাল ফাইবার ব্যবহারের যৌক্তিকতা বিশ্লেষণ করো।
\end{enumerate}
\par

\par\Needspace{7\baselineskip}
\textbf{প্রয়োগ 5: বিষয় ২.5}\par
\par
একটি শিক্ষাপ্রতিষ্ঠান ক্যাম্পাসে তারবিহীন ইন্টারনেট সেবা দিতে চায়। শ্রেণিকক্ষে Wi-Fi থাকবে এবং মিলনায়তন থেকে সরাসরি ভিডিও সম্প্রচার করা হবে। দূরের গ্রামের শিক্ষার্থীদের জন্য স্যাটেলাইটভিত্তিক বিকল্প সংযোগ বিবেচনা করা হচ্ছে। শিক্ষকরা Bluetooth মাইক্রোফোন ব্যবহার করবেন।

\begin{enumerate}
\item ক. তারবিহীন মাধ্যম কী?
\item খ. ব্যক্তিগত স্বল্প দূরত্বের যোগাযোগে Bluetooth উপযোগী কেন?
\item গ. শ্রেণিকক্ষে তারবিহীন ইন্টারনেট ব্যবহারে Wi-Fi উপযোগী কেন? ব্যাখ্যা করো।
\item ঘ. দূরবর্তী গ্রামের শিক্ষার্থীদের জন্য স্যাটেলাইট সংযোগের যৌক্তিকতা, সুবিধা ও সীমাবদ্ধতা বিশ্লেষণ করো।
\end{enumerate}
\par

\par\Needspace{7\baselineskip}
\textbf{প্রয়োগ 6: বিষয় ২.5}\par
স্মার্ট ক্যাম্পাসের নেটওয়ার্কে শিক্ষার্থীরা Wi-Fi দিয়ে ইন্টারনেট এবং শিক্ষকরা Bluetooth ডিভাইস ব্যবহার করবেন। নিরাপত্তা ক্যামেরায় বেতার সংযোগ থাকবে। দূরের শাখা ক্যাম্পাসের সঙ্গে বিকল্প স্যাটেলাইট সংযোগ রাখা হবে।

\begin{enumerate}
\item ক. Wi-Fi কী?
\item খ. বেতার যোগাযোগে নিরাপত্তা গুরুত্বপূর্ণ কেন?
\item গ. শিক্ষার্থীদের ইন্টারনেট ব্যবহারে Wi-Fi কীভাবে কাজ করবে? ব্যাখ্যা করো।
\item ঘ. উদ্দীপক অনুযায়ী Bluetooth, Wi-Fi ও স্যাটেলাইট সংযোগের ব্যবহার বিশ্লেষণ করো।
\end{enumerate}
\par

\par\Needspace{7\baselineskip}
\textbf{প্রয়োগ 7: বিষয় ২.6}\par
\par
একজন শিক্ষার্থী গ্রামের বাড়ি থেকে মোবাইল ডেটায় অনলাইন ক্লাস করেন। ফোনে 4G দেখালেও ভিডিও মাঝে মাঝে আটকে যায়। বাসে শহরে যাওয়ার সময় কল সচল রেখেই সেবাদানকারী টাওয়ার পরিবর্তন হয়। বিদেশে তিনি একই SIM দিয়ে অন্য সমর্থিত নেটওয়ার্কের সেবা নিতে পারেন।

\begin{enumerate}
\item ক. মোবাইল যোগাযোগ কী?
\item খ. হ্যান্ডঅফ কেন প্রয়োজন?
\item গ. উদ্দীপকের হ্যান্ডঅফ ও রোমিং শনাক্ত করে ব্যাখ্যা করো।
\item ঘ. ফোনে 4G দেখালেই কি অনলাইন ক্লাসের ভালো মান নিশ্চিত হয়? ব্যান্ডউইথ, বিলম্ব, সংকেতের শক্তি ও নেটওয়ার্কের ভিড়ের আলোকে বিশ্লেষণ করো।
\end{enumerate}
\par

\par\Needspace{7\baselineskip}
\textbf{প্রয়োগ 8: বিষয় ২.6}\par
শহরের মোবাইল নেটওয়ার্ক উন্নত করতে বিভিন্ন স্থানে নতুন টাওয়ার বসানো হচ্ছে। চলমান ব্যবহারকারীর কলের সময় টাওয়ার বদলালেও কল সচল থাকে। শহরে 4G আছে; তবে ব্যবহারকারীর চাপ বেশি এমন স্থানে গতি কমে। স্মার্ট ট্রাফিক ব্যবস্থা ও IoT ডিভাইসের জন্য ভবিষ্যতে 5G সম্প্রসারণের পরিকল্পনা রয়েছে।

\begin{enumerate}
\item ক. BTS কী?
\item খ. সেলুলার নেটওয়ার্কের এলাকাকে সেলে ভাগ করা হয় কেন?
\item গ. উদ্দীপকে হ্যান্ডঅফের ভূমিকা ব্যাখ্যা করো।
\item ঘ. ভিড়ের এলাকায় 4G-এর গতি কমার কারণ এবং স্মার্ট শহরের সেবায় 5G-এর সহায়ক ভূমিকা বিশ্লেষণ করো।
\end{enumerate}
\par

\par\Needspace{7\baselineskip}
\textbf{প্রয়োগ 9: বিষয় ২.7}\par
\par
একটি কলেজে নতুন ডিজিটাল ব্যবস্থা চালু হয়েছে। ল্যাবের ৪০টি কম্পিউটার একই প্রিন্টার ও ইন্টারনেট সংযোগ ব্যবহার করে। অধ্যক্ষের কার্যালয় ও হিসাব শাখা একই শিক্ষার্থী ডেটাবেজ ব্যবহার করে। কলেজের দুটি শাখা একই শহরে সংযুক্ত এবং প্রধান কার্যালয় অন্য জেলায় অবস্থিত।

\begin{enumerate}
\item ক. কম্পিউটার নেটওয়ার্ক কী?
\item খ. রিসোর্স ভাগাভাগি গুরুত্বপূর্ণ কেন?
\item গ. ল্যাবের নেটওয়ার্ক কোন শ্রেণির হতে পারে? ব্যাখ্যা করো।
\item ঘ. একই শহরের শাখা এবং অন্য জেলার প্রধান কার্যালয়ের সংযোগে MAN ও WAN-এর ব্যবহার বিশ্লেষণ করো।
\end{enumerate}
\par

\par\Needspace{7\baselineskip}
\textbf{প্রয়োগ 10: বিষয় ২.7}\par
হাসপাতালের চিকিৎসকদের কক্ষের কম্পিউটারগুলো একই প্রিন্টার ও রোগীর ডেটাবেজ ব্যবহার করবে। অন্য শাখা একই শহরে যুক্ত থাকবে। প্রধান কার্যালয় অন্য দেশে অবস্থিত এবং অনলাইন প্রতিবেদনব্যবস্থা থাকবে। সরবরাহকারীরা নির্দিষ্ট পোর্টাল দিয়ে ওষুধের মজুতের তথ্য হালনাগাদ করতে পারবেন।

\begin{enumerate}
\item ক. LAN কী?
\item খ. ইন্ট্রানেট ব্যক্তিগত নেটওয়ার্ক কেন?
\item গ. চিকিৎসকদের কক্ষের কম্পিউটারের নেটওয়ার্ক কোন শ্রেণির? ব্যাখ্যা করো।
\item ঘ. শাখা, প্রধান কার্যালয় ও সরবরাহকারীর পোর্টালে MAN, WAN ও Extranet-এর ব্যবহার বিশ্লেষণ করো।
\end{enumerate}
\par

\par\Needspace{7\baselineskip}
\textbf{প্রয়োগ 11: বিষয় ২.8}\par
\par
কলেজের ICT ল্যাব আধুনিক করতে ৪০টি কম্পিউটার, ১টি সার্ভার ও ২টি প্রিন্টার প্রয়োজন। ক্যাম্পাসে Wi-Fi, ইন্টারনেট সংযোগ ও নেটওয়ার্ক নিরাপত্তাও দরকার। ISP থেকে একটি ফাইবার সংযোগ এসেছে। ল্যাবের কিছু অংশে সংকেত দুর্বল।

\begin{enumerate}
\item ক. নেটওয়ার্ক ডিভাইস কী?
\item খ. হাবের তুলনায় সুইচ অধিক কার্যকর কেন?
\item গ. ৪০টি কম্পিউটার ও প্রিন্টারগুলো একই LAN-এ যুক্ত করতে কোন ডিভাইস উপযোগী? ব্যাখ্যা করো।
\item ঘ. ইন্টারনেট সংযোগ, ক্যাম্পাসের Wi-Fi, দুর্বল সংকেত ও নিরাপত্তা নিয়ন্ত্রণে কোন নেটওয়ার্ক ডিভাইসগুলো উপযোগী? বিশ্লেষণ করো।
\end{enumerate}
\par

\par\Needspace{7\baselineskip}
\textbf{প্রয়োগ 12: বিষয় ২.8}\par
হাসপাতালের নেটওয়ার্কে রোগী ব্যবস্থাপনার সার্ভার, চিকিৎসকদের কক্ষের ডেস্কটপ কম্পিউটার, Wi-Fi এলাকা, ইন্টারনেট সংযোগ ও নিরাপত্তাব্যবস্থা প্রয়োজন। ভবন বড় হওয়ায় কিছু স্থানে বেতার সংকেত দুর্বল।

\begin{enumerate}
\item ক. NIC কী?
\item খ. রাউটার ভিন্ন নেটওয়ার্ক সংযুক্ত করতে পারে কেন?
\item গ. চিকিৎসকদের ডেস্কটপ কম্পিউটার ও সার্ভার একই LAN-এ যুক্ত করতে কোন ডিভাইস প্রয়োজন? ব্যাখ্যা করো।
\item ঘ. ইন্টারনেট, Wi-Fi-এর আওতা, দুর্বল সংকেত ও নিরাপত্তা নিয়ন্ত্রণে কোন ডিভাইসগুলো উপযোগী? বিশ্লেষণ করো।
\end{enumerate}
\par

\par\Needspace{7\baselineskip}
\textbf{প্রয়োগ 13: বিষয় ২.10}\par
\par
কলেজ অনলাইন ক্লাস, ডিজিটাল নোট ভাগাভাগি, অ্যাসাইনমেন্ট জমা ও ফলাফল ব্যবস্থাপনা চালু করতে চায়। শিক্ষকরা Google Drive-এ নোট দেবেন। শিক্ষার্থীরা অনলাইনে ডকুমেন্ট সম্পাদনা করবে। কর্তৃপক্ষ ফলাফলের ডেটাবেজ নিয়ন্ত্রিত ও নিরাপদ পরিবেশে রাখতে চায়।

\begin{enumerate}
\item ক. ক্লাউড কম্পিউটিং কী?
\item খ. শিক্ষার্থীদের জন্য ক্লাউড স্টোরেজ সুবিধাজনক কেন?
\item গ. অনলাইনে ডকুমেন্ট সম্পাদনায় কোন ক্লাউড সেবা মডেল ব্যবহৃত হচ্ছে? ব্যাখ্যা করো।
\item ঘ. কলেজের নোট ভাগাভাগি, অনলাইন ক্লাস ও নিরাপদ ফলাফল ব্যবস্থাপনায় Public, Private বা Hybrid Cloud কীভাবে ব্যবহার করা যায়? বিশ্লেষণ করো।
\end{enumerate}
\par

\par\Needspace{7\baselineskip}
\textbf{প্রয়োগ 14: বিষয় ২.10}\par
হাসপাতাল রোগীর প্রতিবেদন, সাক্ষাতের সময় নির্ধারণ, চিকিৎসকের অনলাইন পরামর্শ ও জরুরি ব্যাকআপ ক্লাউডে নিতে চায়। রোগীর তথ্য সংবেদনশীল হওয়ায় কর্তৃপক্ষ নিরাপত্তায় বেশি নিয়ন্ত্রণ রাখতে চায়। অনলাইন পরামর্শ একসঙ্গে অনেক ব্যবহারকারী নেবেন। তাই প্রয়োজন অনুযায়ী সক্ষমতা বাড়ানোর সুবিধা দরকার।

\begin{enumerate}
\item ক. ক্লাউড স্টোরেজ কী?
\item খ. সাধারণ ব্যবহারকারীর জন্য SaaS সহজ কেন?
\item গ. চিকিৎসকদের অনলাইন পরামর্শের প্ল্যাটফর্ম কোন ক্লাউড মডেল ব্যবহার করতে পারে? ব্যাখ্যা করো।
\item ঘ. সংবেদনশীল রোগীর তথ্য এবং চাহিদামতো সক্ষমতা বাড়ানো যায় এমন পরামর্শব্যবস্থায় Hybrid Cloud-এর যৌক্তিকতা বিশ্লেষণ করো।
\end{enumerate}
\par

